\subsection{Neutrino Energy Resolution}\label{sec:ENERGY_RESOLUTION}
The energy resolution is calculated from the calibrated neutrino energies as the \gls{rms} of the relative error of the reconstructed variable versus the truth. The results are shown in \cref{fig:SIGNAL_RMS} for the 4 different reconstruction algorithms. The energy resolution is evaluated as a function of the true neutrino energy, which is obtained from the simulation truth information. This allows us to assess how well the reconstructed energy matches the true energy across a range of neutrino energies, providing insight into the performance of the reconstruction algorithms.

\begin{figure}[ht]
    \centering
    \includegraphics[width=0.49\linewidth]{Images/Reco/neutrino_energy_resolution_hd_1x2x6_centralapa_marley_official_values_rms_variable_drift_reco_nhits_3_comparison.png}
    \includegraphics[width=0.49\linewidth]{Images/Reco/neutrino_energy_resolution_vd_1x8x14_3view_30deg_nominal_marley_official_values_rms_variable_drift_reco_nhits_3_comparison.png}
    \caption[width=1\linewidth]{Neutrino energy resolution for the different reconstruction algorithms. (left) \gls{hd} geometry. (right) \gls{vd} geometry.\label{fig:SIGNAL_RMS}}
\end{figure}

As expected, the \emph{ClusterEnergy} algorithms has the worst energy resolution for neutrino energies below \SI{10}{\mega\electronvolt}, as it does not take into account the energy deposited by adjacent gamma blips which can amount to more than 50\% of the initial neutrino energy in this range. The \emph{TotalEnergy} algorithm, on the other hand, has a good energy resolution in the full range for \gls{hd}, but it is significantly worse for the \gls{vd} geometry due to the higher background rate and lower clustering efficiency.

The \emph{SelectedEnergy} and \emph{SolarEnergy} algorithms provide a significant improvement in the energy resolution for neutrino energies below \SI{30}{\mega\electronvolt} by maximizing the adjacent cluster significance. The advantages of the \emph{SolarEnergy} algorithm becomes apparent for event candidate reconstruction in the context of solar neutrino analyses, where the signal energy range is expected to be dominated by external gammas and neutrons from the cavern. The latter produce a gamma cascade which results in an enhanced adjacent cluster activity. If not limited algorithmically (e.g. by applying a maximum adjacent cluster energy as in the \emph{SolarEnergy} method), this can lead to a significant overestimation of the event's energy. Further details on this will be provided in the solar neutrino analysis note.

\begin{figure}
    \centering
    \includegraphics[width=.49\linewidth]{Images/Reco/neutrino_energy_resolution_hd_1x2x6_centralapa_marley_official_values_rms_solarenergy_drift_nhits_3_scan.png}
    \includegraphics[width=.49\linewidth]{Images/Reco/neutrino_energy_resolution_vd_1x8x14_3view_30deg_nominal_marley_official_values_rms_solarenergy_drift_nhits_3_scan.png}
    \caption{(left) \gls{hd} (right) \gls{vd}. Neutrino energy resolution with and without flash-matching. The labels \textbf{True}, \textbf{Reco} and \textbf{None} represent the cases for which the drift coordinate is taken from truth information, from the flash-matching step and with an average drift correction, respectively.\label
    {fig:FLASH_MATCHING_EFFECT}}
\end{figure}

\Cref{fig:FLASH_MATCHING_EFFECT} shows the effect of the flash-matching step on the energy resolution. The flash-matching step allows to associate a \gls{tpc} cluster with the corresponding \gls{pds} flash and hence obtain the correct drift coordinate of the interaction. The difference between the flash-matched and non-flash-matched cases are small for \gls{hd}, indicating that the inefficiencies in the energy reconstruction are dominated by other factors, such as clustering inefficiencies and intrinsic energy losses due to nuclear effects and undetected particles. In the case of \gls{vd}, the flash-matching step has a more noticeable effect on the energy resolution (\SI{\sim5}{\percent} in the \qtyrange{4}{20}{\mega\electronvolt} range). This is expected due to the larger drift volume.

\begin{figure}[ht!]
    \centering
    \includegraphics[width=0.49\linewidth]{Images/Reco/neutrino_energy_resolution_hd_1x2x6_vd_1x8x14_3view_30deg_marley_official_values_rms_solarenergy_drift_reco_nhits_3_comparison.png}
    \includegraphics[width=0.49\linewidth]{Images/Reco/neutrino_energy_resolution_marley_official_values_rms_solarenergy_drift_reco_nhits_3_comparison.png}
    \caption[width=1\linewidth]{Energy resolution comparison between all geometry and background configurations. (left) Comparison between the \seqsplit{hd\_1x2x6\_centralAPA} and \seqsplit{vd\_1x8x14\_3view\_30deg} geometries. (right) Comparison between the different background scenarios for the \seqsplit{vd\_1x8x14\_3view\_30deg} geometry.\label{fig:ENERGY_RESOLUTION_COMPARISON}}
\end{figure}

Finnaly, \cref{fig:ENERGY_RESOLUTION_COMPARISON} shows the energy resolution comparison between all geometry and background configurations. We can see that the energy resolution is similar between the \seqsplit{hd\_1x2x6} and \seqsplit{vd\_1x8x14\_3view\_30deg} geometries, with a noticeable advantage for the former in the configurations with background (\seqsplit{hd\_1x2x6\_centralAPA} \& \seqsplit{hd\_1x2x6\_lateralAPA} vs \seqsplit{vd\_1x8x14\_3view\_30deg\_nominal} \& \seqsplit{vd\_1x8x14\_3view\_30deg\_shielded}).